\chapter{Conclusion}

This thesis presented the complete design, development, and validation of an autonomous robotic 
manipulator tailored for the selective harvesting of green asparagus. The primary objective was 
to bridge the gap between theoretical robotics and practical agricultural application by 
creating a system capable of identifying, locating, and harvesting asparagus spears in a 
simulated environment. The work encompassed the complete mechatronic design of a custom 
2-degree-of-freedom (DOF) robotic arm, and a separate specialized cutting and gripping 
end-effector. On the software front, a robust control architecture was established, 
integrating high-level navigation logic written in Python with low-level firmware executing 
on distributed microcontrollers. A sophisticated computer vision pipeline was developed, 
utilizing a Mask R-CNN deep learning model for instance segmentation and a 3D point cloud 
processing sequence for precise spatial localization. The integration of 2D visual data with 
3D depth information, refined through a centroid-filtering technique, enabled the reliable 
calculation of harvesting targets. Experimental validation demonstrated the system's 
effectiveness, with the kinematic model verified through rigorous pose calibration and the 
end-to-end harvesting cycle successfully executed. The system proved capable of autonomously 
detecting target spears, planning collision-free trajectories, and performing the delicate 
cut-and-hold operation required for asparagus harvesting.

While the developed prototype successfully demonstrated the proof-of-concept for autonomous 
harvesting, several limitations remain. The assessment of positioning accuracy and harvesting 
success relied primarily on visual estimation due to the lack of external metrology equipment, 
such as laser trackers, which prevented a quantitative analysis of sub-millimeter kinematic 
errors. Furthermore, the experiments were conducted in a controlled laboratory setting with 
consistent lighting conditions and stable target positioning. Consequently, the system has 
not yet been stress-tested against the unpredictable variables found in an actual agricultural 
field, such as variable sunlight, wind, or occlusions by foliage. Additionally, the system 
currently operates from a fixed base, whereas in a real-world scenario, the manipulator must 
be mounted on a mobile platform to traverse crop rows, introducing new challenges in dynamic 
localization and vibration stabilization.

To evolve this prototype into a field-ready agricultural product, future research should focus 
on conducting extensive testing in real asparagus fields to evaluate the computer vision 
system's robustness against varying lighting conditions and complex background clutter. 
Essential to this evolution is the integration of the manipulator onto an autonomous ground 
vehicle (rover) to enable continuous harvesting along crop rows, coordinating the arm's motion 
with the vehicle's odometry. Further improvements could include implementing a closed-loop 
visual servoing system to provide real-time feedback during the reaching phase, thereby 
correcting for any target movement or calibration drift. Finally, optimizing the harvesting 
cycle by executing vision processing and arm motion simultaneously would significantly reduce 
the time per harvest, increasing the commercial viability of the system.

